\begin{textsample}{POS Dim 1 – sc – Score 57.00 – sc/sc\_em\_73.txt}  \label{ex:f1_pos_002}
5.5.4 Organização das unidades curriculares da trilha de aprofundamento Unidade Curricular I : Minha morada, tua morada, nossa morada Nesta unidade são apresentados objetos de conhecimento \textbf{que} possibilitam ao/à estudante compreender -se como \textbf{sujeito} histórico em seu lugar de vivência.
Almeja -se \textbf{uma} \textbf{abordagem} de conteúdos \textbf{que} pressupõem o reconhecimento e o respeito das diferenças pessoais, grupais ou culturais sem reduzi -las à compreensão de \textbf{um} eu.
Problematizam -se processos de homogeneização das diferentes identidades, compreendendo, assim, \textbf{que} as relações socioculturais constituem os \textbf{sujeitos} históricos nas organizações de suas vidas sociais e políticas, nas suas relações com o ambiente e com outros grupos, na produção e reprodução de suas existências.
Nesta unidade, ganham destaque as construções, as moradas e as moradias, o patrimônio histórico alimentar, cultural, artístico e religioso, \textbf{que} devem ser tratados pela escola como saberes necessários à compreensão das identidades das juventudes catarinenses.
Unidade Curricular II : Discursos públicos e privados nas esferas da justiça social e da diversidade como direito humano Nesta unidade, são apresentados objetos de conhecimento \textbf{que} visam a contribuir para a inserção do(a) estudante no exercício da cidadania, percebendo esta ação como participação ativa, ou seja, desnaturalizada, consciente, responsável e transformadora do contexto local e global.
Espera -se, assim, \textbf{que} o/a estudante se torne capaz de identificar as violências ( simbólicas e não simbólicas ) e os estereótipos presentes em diversos discursos públicos e/ou privados, adotando \textbf{uma} postura de repulsa frente a discursos racistas, xenofóbicos e tantos outros preconceitos.
No esteio das atividades, destacamos a importância de discussões e conhecimentos acerca das legislações resultantes do engajamento dos indígenas, das comunidades de origem afro, de quilombolas, em âmbito federal, estadual e municipal.
Além disso, \textbf{que} sejam evidenciadas produções artístico-literárias, na promoção da visibilidade de autores(as) locais, regionais e estaduais.
Unidade Curricular III : Povos originários e conflitos por terra e território em Santa Catarina : guarani, kaingang e xokleng/laklanõ Nesta unidade, são apresentados objetos de conhecimento \textbf{que} enfocam os povos originários guarani, kaingang e xokleng/laklanõ, \textbf{historicamente} dizimados e \textbf{que} permanecem invisibilizados ou tomados como figuras folclóricas do passado, desenvolvendo -se nos(as) estudantes \textbf{uma} perspectiva de positivação das identidades indígenas, não somente como \textbf{uma} outra cultura a fazer parte de \textbf{uma} “ mescla transitória ”.
Ambiciona -se \textbf{que} conheçam os modos de esses povos \textbf{verem} e serem no mundo, \textbf{que} conheçam a existência de \textbf{uma} outra ética, de \textbf{uma} outra moral e de formas diferentes de resolução dos problemas e conflitos cotidianos, \textbf{constitutivos} do capital cultural desses povos de matriz filosófica ameríndia.
Mais do \textbf{que} isso, \textbf{que} problematizem toda e qualquer manifestação eurocêntrica na produção e escolha dos conhecimentos em todas as áreas.
Isso \textbf{diz} respeito aos processos de hierarquização \textbf{que} se estabelecem entre o \textbf{que} é o “ normal ” e o \textbf{que} é o “ diverso ” nas escolhas \textbf{metodológicas}, bem como no material didático, em \textbf{que} a representação amazônica, silvícola e quinhentista obscurece a presença indígena em território catarinense e em outros territórios brasileiros.
Unidade Curricular IV : Reconhecimento e respeito às identidades e diversidades humanas e suas \textbf{territorialidades} : comunidades originárias de povos indígenas, comunidades de origem afro, afro-brasileira, quilombolas, caboclos, migrantes/emigrantes, entre outros, em Santa Catarina Nesta unidade, são apresentados objetos de conhecimento \textbf{que} levam o/a estudante a compreender \textbf{que} existem diversas formas de perceber o mundo, \textbf{remetendo} às diferentes identidades \textbf{constitutivas} dos seres humanos, de suas organizações sociais, etnias, nacionalidades, religiosidades, entre outras.
Visa, assim, a compreender a pluralidade e a diversidade étnicas presentes no território catarinense, identificando formas plurais de organização social e manifestações culturais existentes, desde a chegada de outros povos a este estado.
Defende -se, sobretudo, o respeito e a construção de relações recíprocas por meio do diálogo num movimento de aproximação, de conhecimento mútuo, de interação e coexistência entre as diferentes culturas.
As \textbf{abordagens} devem propiciar o conhecimento e a reflexão não somente de festejos, músicas, danças, jogos oficializados e produções artístico-literárias do cânone, mas, principalmente, a inclusão de manifestações marginalizadas, protagonizadas por indígenas, africanos e afro-brasileiros, em \textbf{um} movimento de escuta de quem ali se expressa e o \textbf{que} expressa.
Sugere -se \textbf{que} se façam estudos acerca da \textbf{influência} africana, indígena e do imigrante no vocabulário, na morfologia, na sintaxe, na fonologia e pronúncia, valorizando a diversidade linguística.
Quadro 56 — Minha morada, tua morada, nossa morada Unidade Curricular I — Minha morada, tua morada, nossa morada Carga horária : 35 h ou 55 h, a \textbf{depender} da matriz curricular em funcionamento na unidade escolar Habilidades da trilha associadas aos eixos estruturantes Investigação científica - Analisar, elaborar hipóteses, formular e resolver problemas resultantes de investigações científicas, utilizando conhecimentos das diversas áreas de conhecimento, para criar e/ou propor soluções para problemas sociais diversos.
Processos criativos - Reconhecer como \textbf{legítimas} as diferentes manifestações criativas, artísticas e culturais de distintos grupos étnicos, por meio de vivências presenciais e virtuais para ampliar a visão de mundo, sensibilidade, criticidade e criatividade. - Refletir sobre seu próprio desenvolvimento e sobre seus objetivos presentes e futuros, reconhecendo e respeitando as diferentes aspirações de vida, bem como as escolhas, esforços e ações em relação à sua vida pessoal e cidadã.
Mediação e intervenção sociocultural - Promover capacidade de identificar as atitudes e padrões de comportamento \textbf{que} perpetuam o racismo e a xenofobia.
Desta forma, atuar como agentes na promoção do exercício da alteridade, diante das formas sociais de \textbf{exclusão} dos diversos grupos étnicos, e assim, tomar decisões conscientes, colaborativas e responsáveis.
Habilidades das trilhas de aprofundamento vinculadas às áreas do conhecimento - Identificar e \textbf{explicitar} as questões sobre o ser humano como \textbf{sujeito} individual e coletivo, constituído por diferentes culturas, linguagens, visões de mundo, etnias e demais interseccionalidades. - Mobilizar conhecimentos e recursos das práticas das linguagens propondo ações para o conhecimento de si e conhecimento do outro na corporeidade e subjetividades \textbf{que} compõem os diferentes \textbf{sujeitos} e seus modos de vida. - Analisar questões sociais e culturais (des)construindo o \textbf{que} está posto em \textbf{um} movimento de acolher, incorporar e \textbf{agregar} valores importantes para si e para o outro. - Propor ações de reconhecimento da alteridade nos lugares de vivência. - Compreender os diferentes sentimentos de pertencimento e de relação com a terra, \textbf{que} prezam pela relação de cuidado, sustentabilidade e de respeito ao tempo da natureza. - Analisar a \textbf{influência} dos fluxos econômicos e populacionais na formação socioeconômica e territorial de Santa Catarina, compreendendo os conflitos e as tensões históricas, em diferentes temporalidades e espacialidades. - Compreender os lugares de vivência e em suas histórias familiares e/ou da comunidade, elementos de distintas culturas. - Identificar e comparar aspectos culturais dos grupos sociais de seus lugares de vivência, considerando o campo e a cidade. - Identificar as características socioculturais observadas nos lugares de vivências e suas transformações ao longo do tempo histórico. - Identificar a arquitetura e formas de moradias existentes nos lugares de vivência, percebendo -a em suas diferenças, diversidades e também desigualdades.
Objetos de conhecimento vinculados às habilidades das áreas do conhecimento - Identidades individual e coletiva - Conhecimento de si e do outro, corporeidade e subjetividades - Reconhecimento e alteridade - Sentimento de pertencimento e relação com a terra - Dinâmicas populacionais e formação socioeconômica do território catarinense - Lugares de vivência, histórias familiares e da comunidade - Características socioculturais nos lugares de vivência - Moradias catarinenses e sua relação com a cultura, tradição e contextos econômicos Unidade Curricular II — Discursos públicos e privados nas esferas da justiça social e da diversidade como direito humano Carga horária : 35 h ou 55 h, a \textbf{depender} da matriz curricular em funcionamento na unidade escolar Habilidades da trilha associadas aos eixos estruturantes Investigação científica - Analisar, elaborar hipóteses, formular e resolver problemas resultantes de investigações científicas, utilizando conhecimentos das diversas áreas de conhecimento, para criar e/ou propor soluções para problemas sociais diversos.
Processos criativos - Reconhecer como \textbf{legítimas} as diferentes manifestações criativas, artísticas e culturais de distintos grupos étnicos, por meio de vivências presenciais e virtuais para ampliar a visão de mundo, sensibilidade, criticidade e criatividade. - Refletir sobre seu próprio desenvolvimento e sobre seus objetivos presentes e futuros, reconhecendo e respeitando as diferentes aspirações de vida, bem como as escolhas, esforços e ações em relação à sua vida pessoal e cidadã.
Mediação e intervenção sociocultural - Promover capacidade de identificar as atitudes e padrões de comportamento \textbf{que} perpetuam o racismo e a xenofobia.
Desta forma, atuar como agentes na promoção do exercício da alteridade, diante das formas sociais de \textbf{exclusão} dos diversos grupos étnicos, e assim, tomar decisões conscientes, colaborativas e responsáveis.
Habilidades das trilhas de aprofundamento vinculadas às áreas do conhecimento - Desenvolver projetos utilizando as práticas de linguagens para formular propostas articuladas, \textbf{que} evidenciem a cidadania consciente e responsável. - Avaliar o exercício da cidadania como oportunidades de crescimento individual e coletivo quando exercido com consciência e responsabilidade, \textbf{constituintes} indispensáveis da vida ética. - Propor questões \textbf{que} problematizem a cidadania na participação ativa no contexto de vivência. - Valorizar os processos \textbf{que} contribuem para a formação de \textbf{sujeitos} éticos, com vista ao respeito às práticas de liberdade, cooperação, convivência democrática, solidariedade e fraternidade. - Identificar nas diferentes línguas e linguagens ( imagéticas, gestuais e multimodais ) estereótipos e propagação do racismo, da xenofobia e de violência contra gênero, etnia, geração, idade, classe, \textbf{raça}, idosos, juventudes, lesionados/deficientes, entre outros. - Promover momentos de reflexão e debate sobre a importância do respeito das crenças e posturas, políticas, religiosas e espirituais, evitando atitudes e pensamentos extremistas. - Compreender as identidades culturais na sua complexidade, considerando -as como \textbf{um} processo em construção, múltiplo, não fixo e/ou \textbf{central}.
Objetos de conhecimento vinculados às habilidades das áreas do conhecimento - Práticas de linguagens em cidadania - Exercício da cidadania e ética - Linguagens, estereótipos e preconceito - Decolonialidade - Respeito mútuo, justiça e solidariedade - Políticas públicas : participação direta e indireta do cidadão - Direitos \textbf{humanos} e direitos individuais e sociais Unidade Curricular III — Povos originários e conflitos por terra e território em Santa Catarina : guarani, kaingang e xokleng/laklanõ Carga horária : 35 h ou 55 h Habilidades da trilha associadas aos eixos estruturantes Investigação científica - Investigar, analisar, elaborar e testar hipóteses, formular e resolver problemas, criando soluções a partir dos conhecimentos de diferentes áreas. - Utilizar informações, conhecimentos e ideias resultantes de investigações científicas para criar e/ou propor soluções para problemas sociais diversos.
Processos criativos - Reconhecer como \textbf{legítimas} as diferentes manifestações criativas, artísticas e culturais de distintos grupos étnicos, por meio de vivências presenciais e virtuais para ampliar a visão de mundo, sensibilidade, criticidade e criatividade. - Promover capacidade de identificar as atitudes e padrões de comportamento \textbf{que} perpetuam o racismo e a xenofobia.
Desta forma, atuar como agentes na promoção do exercício da alteridade, diante das formas sociais de \textbf{exclusão} dos diversos grupos étnicos, e assim, tomar decisões conscientes, colaborativas e responsáveis.
Mediação e intervenção sociocultural - Refletir sobre seu próprio desenvolvimento e sobre seus objetivos presentes e futuros, reconhecendo e respeitando as diferentes aspirações de vida, bem como as escolhas, esforços e ações em relação à sua vida pessoal, profissional e cidadã.
Empreendedorismo - Perceber as diferentes formas de empreender, seja no campo do bem-viver social, da cultura, das artes, da economia, bem como das questões espirituais e religiosas, desvinculando tal concepção exclusivamente ao mercado de trabalho.
Habilidades das trilhas de aprofundamento vinculadas às áreas do conhecimento - Conhecer e compreender a importância da História e Cultura Indígena em Santa Catarina, bem como a trajetória dos povos Guarani, Kaingang e Xokleng/Laklanõ até o momento presente, \textbf{que} \textbf{historicamente} foram invisibilizadas, ou tomadas como figuras folclóricas do passado. - Compreender \textbf{que} existe \textbf{uma} outra ética, \textbf{uma} outra moral e formas diferentes de resolução dos problemas e conflitos cotidianos, \textbf{constitutivos} do capital cultural dos povos de matriz filosófica ameríndia. - Historicizar as próprias condições em \textbf{que} se elaboraram as representações dos \textbf{sujeitos} indígenas em Santa Catarina diante da presença de imigrantes europeus, e \textbf{que} também \textbf{dizem} respeito aos processos históricos \textbf{que} nos constituíram como nação. - Desenvolver \textbf{uma} perspectiva de positivação das identidades indígenas e não somente como \textbf{uma} outra cultura a fazer parte de \textbf{uma} “ mescla transitória ”. - Problematizar toda e qualquer manifestação de processos hierárquicos \textbf{que} estabelecem \textbf{um} “ normal ”, obscurecendo a presença dos povos originários em território catarinense. - Reconhecer os processos de empreendedorismo dos \textbf{sujeitos} indígenas em Santa Catarina na perspectiva do etnodesenvolvimento. - Perceber a dimensão dos princípios da reciprocidade, da fraternidade entre as pessoas, da convivência e respeito para com a natureza exercida pelos povos indígenas — o Bem Viver — como possibilidade de escolha para a vida nos espaços urbanos, assegurando a existência humana. - Reconhecer a história das culturas indígenas, \textbf{constitutivas} da diversidade e para o exercício da cidadania, considerando as \textbf{territorialidades}. - Promover ações de combate às atitudes e situações fomentadoras de todo tipo de discriminação e injustiça social, em especial contra os \textbf{sujeitos} indígenas. - Reconhecer e respeitar a presença das culturas indígenas no campo do conhecimento científico, bem como das demais produções de conhecimento. - Reconhecer produtos e/ou processos criativos por meio de fruição, vivências sobre a arte, arte \textbf{literária} e cultura dos povos indígenas. - Refletir sobre os debates acerca do discurso e das representações da conquista da América e as formas de organização dos indígenas e europeus, e das resistências indígenas considerando \textbf{uma} releitura apresentada pela perspectiva decolonial.
Objetos de conhecimento vinculados às habilidades das áreas do conhecimento - História e Cultura Indígena em Santa Catarina : guarani, kaingang e xokleng/laklanõ - Ética, moral e filosofia ameríndia - Representações dos \textbf{sujeitos} indígenas - Identidades indígenas - Processos hierárquicos e reconhecimento - Etnodesenvolvimento - Saberes, fazeres e conhecimentos indígenas - História das culturas indígenas, diversidades e cidadania - Combate à discriminação e à injustiça social - Cientificidade nos saberes e conhecimentos dos povos originários indígenas - Desconstrução do modelo eurocêntrico da Ciência e inclusão dos saberes científicos indígenas Unidade Curricular IV — Reconhecimento e respeito às identidades e diversidades humanas e suas \textbf{territorialidades} : comunidades originárias de povos indígenas, comunidades de origem afro, quilombolas, os caboclos, migrantes/emigrantes, entre outros, em Santa Catarina Carga horária : 55 h ou 75 h, a \textbf{depender} da matriz curricular em funcionamento na unidade escolar Habilidades da trilha associadas aos eixos estruturantes Investigação científica - Investigar, analisar, elaborar e testar hipóteses, formular e resolver problemas, criando soluções a partir dos conhecimentos de diferentes áreas. - Utilizar informações, conhecimentos e ideias resultantes de investigações científicas para criar e/ou propor soluções para problemas sociais diversos.
Processos criativos - Reconhecer como \textbf{legítimas} as diferentes manifestações criativas, artísticas e culturais de distintos grupos étnicos, por meio de vivências presenciais e virtuais para ampliar a visão de mundo, sensibilidade, criticidade e criatividade. - Promover capacidade de identificar as atitudes e padrões de comportamento \textbf{que} perpetuam o racismo e a xenofobia.
Desta forma, atuar como agentes na promoção do exercício da alteridade, diante das formas sociais de \textbf{exclusão} dos diversos grupos étnicos, e assim, tomar decisões conscientes, colaborativas e responsáveis.
Mediação e intervenção sociocultural - Refletir sobre seu próprio desenvolvimento e sobre seus objetivos presentes e futuros, reconhecendo e respeitando as diferentes aspirações de vida, bem como as escolhas, esforços e ações em relação à sua vida pessoal, profissional e cidadã.
Empreendedorismo - Perceber as diferentes formas de empreender, seja no campo do bem-viver social, da cultura, das artes, da economia, bem como das questões espirituais e religiosas, desvinculando tal concepção exclusivamente ao mercado de trabalho.
Habilidades das trilhas de aprofundamento vinculadas às áreas do conhecimento - Compreender a diversidade como característica dos indivíduos, percebendo como diversos em suas experiências de vida históricas e culturais e em suas identidades, considerando organizações sociais, etnias, nacionalidades, religiosidades, entre outras. - Entender \textbf{que} as relações socioculturais constituem os \textbf{sujeitos} históricos, nas organizações de suas vidas sociais e políticas, nas suas relações com o ambiente e com outros grupos, na produção e reprodução de suas existências. - Compreender a pluralidade e a diversidade no território catarinense, identificando as plurais formas de \textbf{percepção} da realidade, da organização social e manifestações culturais e seus diferentes \textbf{sujeitos} históricos. - Investigar, \textbf{sistematizar} e descrever, por meio de múltiplas linguagens ( corporal, oral, escrita, audiovisual ), as brincadeiras e os jogos e festas \textbf{populares} de matriz indígena e africanas e dos imigrantes e migrantes, dos caboclos, quilombolas, reconhecendo a importância desse patrimônio \textbf{histórico-cultural} na preservação das diferentes culturas. - Reconhecer as diferentes formas de empreendedorismo dos povos catarinenses. - Perceber as diferentes relações \textbf{que} se estabelecem com os territórios e as narrativas sobre estes. - Identificar, reconhecer e compreender as identidades no âmbito das diversidades das matrizes histórica e culturais catarinenses ( povos originários, afro-brasileiros, quilombolas, entre outros ), considerando o multiculturalismo nas suas historicidades e \textbf{territorialidades}. - Compreender os Direitos Humanos para a formação intercultural, percebendo as \textbf{singularidades} e diversidade dos \textbf{sujeitos} históricos e grupos sociais, inclusive por meio das perspectivas decolonizantes, nas diversas escalas. - Conhecer e valorizar o patrimônio cultural, material e imaterial, das culturas no território catarinense, incluindo as diversas matrizes indígenas, africanas e europeias, de diferentes épocas, favorecendo a construção de vocabulário e repertório relativos às diferentes linguagens artísticas. - Compreender a disputa territorial do Contestado em Santa Catarina com seus \textbf{desdobramentos} pós movimento do Contestado ( territorial, social, econômico, político ). - Reconhecer e valorizar os diferentes grupos étnicos para \textbf{uma} educação de combate ao racismo e à xenofobia e refletir sobre o impacto da ciência na vida social, nas relações étnico-raciais. - Identificar e descrever territórios étnico-culturais existentes no Brasil, tais como terras indígenas e de comunidades remanescentes de quilombos, reconhecendo a legitimidade da demarcação desses territórios. - Reconhecer, valorizar e respeitar as manifestações da cultura corporal e de movimento das diferentes etnias do território catarinense. - Reconhecer as línguas e seu importante papel na constituição da identidade, os processos de cooficialização dessas línguas e os discursos e práticas preconceituosas e discriminatórias em relação a estas. - Compreender as desigualdades sociais e a luta pela terra dos indígenas, dos quilombolas e dos camponeses, no território catarinense. - Reconhecer as \textbf{territorialidades} dos povos indígenas, das comunidades remanescentes de quilombos, povos das florestas e do cerrado, de \textbf{ribeirinhos} e caiçaras, do campo e da cidade, como direitos legais dessas comunidades. - Compreender a desconstrução da história catarinense eurocentrada e a adoção de \textbf{um} viés decolonial com os estudos dos quilombos, e das tradições africanas e afrocatarinenses. - Conhecer e analisar a diversidade do patrimônio etnocultural e reconhecer a diversidade cultural como \textbf{um} direito dos povos e dos indivíduos e elemento como fortalecimento da democracia.
Objetos de conhecimento vinculados às habilidades das áreas do conhecimento - Diversidade e identidades - Identidades e cultura - Pluralidade e diversidade no território catarinense - Jogos, brincadeiras, arte e festividades \textbf{populares} como patrimônio \textbf{histórico-cultural} - Empreendedorismo - Território e narrativa - Desigualdades e racismo estrutural - Identidades e diversidades catarinenses - Direitos Humanos e formação intercultural - História da Arte de matriz africana, afro-brasileira/catarinense e indígena - Contestado e território - Combate ao racismo e à xenofobia - Quilombos e terras indígenas - Manifestações culturais e corpo - Diversidade linguística : plurilinguismo em território catarinense - Luta pela terra - Terra como direito - Eurocentrismo e história de Santa Catarina - Diversidade etnocultural Fonte : Elaboração dos \textbf{autores} ( 2020 ). 5.5.6 Orientações \textbf{metodológicas} A sociedade está sempre em transformação, principalmente nas diferentes formas de comunicação, relacionamentos e organização espacial, e essas mudanças se refletem na escola, principalmente pelas diferentes formas com \textbf{que} estamos nos relacionando, comunicando e organizando espacialmente.
Os(as) estudantes do ensino médio também estão experimentando esta realidade, principalmente com o acesso às Tecnologias Digitais de Informação e Comunicação ( TDICs ), \textbf{que} permitem a conexão com diversas informações e levam a assumir \textbf{posicionamentos} em vários contextos na sociedade.
No processo educacional, as TDICs, a realidade do/a estudante, da comunidade escolar e das diversas dimensões escalares das ações da sociedade estão \textbf{demandando} \textbf{uma} aprendizagem ampla e aprofundada, conforme os avanços das pesquisas em Psicologia, Neurociência e Educação. \textbf{Uma} das estratégias educacionais \textbf{que} podem atender a esta demanda são as metodologias ativas, fundamentadas em diferentes \textbf{abordagens} \textbf{teóricas}, \textbf{que} têm por objetivo provocar nos(as) estudantes \textbf{um} aprendizado de forma autônoma e participativa, utilizando -se de situações reais, problematizando \textbf{um} estudo de caso, permitindo refletir e tomar iniciativas, sugerindo resolução de problemas e construção de conhecimentos.
Nesta proposta \textbf{metodológica}, o professor se \textbf{configura} como coadjuvante no processo educacional, permitindo aos(às) estudantes o protagonismo de seu aprendizado ( MORAN, 2018 ).
É significativo \textbf{que} a ação docente possibilite ao(à) estudante experiências educativas associadas ao seu contexto em interação com outros aspectos sociais e culturais, e a realidade \textbf{contemporânea}, na promoção de sua formação pessoal, profissional e cidadã, bem como no trabalho relacionado a valores como ética, liberdade, democracia, justiça social, pluralidade, solidariedade e sustentabilidade ( \textbf{SANTA} CATARINA, 2019 ).
A trilha integrada “ Atelier do Território Catarinense : Identidades, Pluralidades e Diversidades ” requer estratégias de ensino e aprendizagem \textbf{que} priorizem sequências didáticas \textbf{que} proporcionem este protagonismo estudantil, \textbf{permeadas} por eixos estruturantes \textbf{que} \textbf{consigam} problematizar, contextualizar, pesquisar e propor resolução de problemas, mediante processos de observação, comparação, investigação, análise, interpretação, reflexão, síntese, elaboração de textos argumentativos, informativos e crônicas, entre outros.
Ou seja, ações \textbf{que} favoreçam aos(às) estudantes o ressignificar conhecimentos e sinalizar possíveis alternativas de respostas aos dilemas presentes no seu cotidiano ( \textbf{SANTA} CATARINA, 2014 ).
O protagonismo estudantil se destaca a partir da cidadania ativa territorial, \textbf{aproximando} o poder local ( câmaras, poderes legislativos, prefeituras, executivos e fóruns, poderes judiciários, entre outros ), por \textbf{intermédio} da participação dos jovens e de suas escolas.
Destaca -se ao \textbf{fomentar} ações \textbf{que} contribuam para \textbf{um} desenvolvimento sustentável das localidades e dos municípios ; a valorizar o estudo de caso como trabalho experimental sobre problemas locais, assim como a promover simulações, jogos, brincadeiras, experiências, vivências, utilização de tecnologias de informação.
Vale \textbf{salientar} \textbf{que} “ jogos e brincadeiras não se reduzem a estratégias de aprendizagem, porque são parte integrante da dimensão humana e do processo de elaboração do pensamento.
Ampliam -se, assim, as possibilidades de intervenção, questionamento e transformação da realidade.
Ao longo desse processo, a \textbf{sistematização}, por meio de registros, relatos e outros instrumentos, é subsídio para a identificação da apropriação de conhecimentos ” ( \textbf{SANTA} CATARINA, 2014, p. 149 ).
Destacamos três propostas \textbf{metodológicas} com base na metodologia ativa, \textbf{que} são \textbf{inerentes} ao protagonismo estudantil.
Elas permitirão \textbf{uma} dinamização das aulas, o incentivo ao pensamento complexo e científico, o \textbf{estímulo} à criatividade, a cooperação em todas as atividades e a demonstração de \textbf{que} todos(as) estudantes são constituídos de saberes ( \textbf{FREIRE}, 1996 ). a ) Modelo das Nações Unidas ( Modelo ONU ) O Modelo das Nações Unidas é \textbf{uma} prática estudantil através de simulações sobre o funcionamento da ONU e a dinâmica das relações diplomáticas internacionais, assumindo os papéis de líderes mundiais e tomadores de decisões internacionais.
De acordo com Randig : > As simulações acadêmicas de conferências e organismos multilaterais já ocorriam antes mesmo da criação da ONU.
Na década de 1920, Harvard já organizava, entre seus alunos, simulações da Sociedade das Nações.
Desde 1955, a Universidade sedia ininterruptamente edições anuais do Harvard National Model United Nations ( HNMUN ), \textbf{que} reúne cerca de dois mil estudantes universitários a cada edição.
A instituição organiza também, desde 1992, o World Model United Nations ( WorldMUN ), cujas conferências são a cada ano sediadas em \textbf{um} diferente país.
Cambridge, Puebla, Taipé, Pequim, Sharm El-Sheikh, Atenas, Heidelberg e Belo Horizonte são algumas das cidades \textbf{que} já sediaram o WorldMUN, recebendo, por \textbf{uma} semana, cerca de dois mil estudantes de todo o mundo.
No Brasil, a primeira conferência de simulação acadêmica das Nações Unidas foi o Amun – o Americas Model United Nations – fundado em 1998 e organizado de forma voluntária por estudantes do curso de graduação em Relações Internacionais da Universidade de Brasília ( UnB ).
A iniciativa \textbf{despertou} enorme interesse desde sua primeira edição, trazendo a Brasília estudantes universitários de todo o Brasil e de mais de \textbf{uma} dezena de outros países.
Desde então, a cultura das “ simulações ” espalhou -se por todo o País : \textbf{hoje}, são dezenas as conferências – tanto universitárias quanto para secundaristas – \textbf{que} ocorrem anualmente nas cinco regiões do Brasil.
Muitas instituições de ensino médio e superior possuem seus próprios clubes e associações permanentes de simulação e, em diversas escolas brasileiras – sobretudo em Brasília –, simulações das Nações Unidas tornaram -se atividade constante e obrigatória para todos os alunos ( 2010, p. 62 ).
Conforme esse mesmo \textbf{autor} e o Harvard Model United Nations ( 2020 ), as simulações do Modelo ONU também podem ser praticadas na educação básica na etapa do ensino médio, consagrando o protagonismo estudantil, por estarem os(as) estudantes diante de diferentes situações políticas, sociais, ambientais, econômicas e culturais nas diversas escalas do espaço geográfico, permitindo problematizarem, elaborarem hipóteses e traçar objetivos, conforme as evidências.
Tais situações permitem apropriar -se de conceitos, conhecimentos e habilidades para a compreensão das relações de poder e a se posicionar de forma crítica diante dos fatos.
Esta atividade estimula e promove a empatia e o respeito pela diversidade humana e ambiental, além de oportunizar debates sobre questões \textbf{urgentes} do cotidiano e esboçar soluções inovadoras e criativas.
Os participantes poderão desenvolver outras habilidades ao longo desta prática, como a oratória, a negociação, o trabalho em equipe, a liderança e a elaboração de políticas.
Na trilha “ Atelier do Território Catarinense : Identidades, Pluralidades e Diversidades ”, sugere -se trazer o Modelo ONU para a escala local e as demais, conforme as demandas, e promover o desenvolvimento do pensamento crítico sobre a realidade do entorno e \textbf{que} os/as estudantes pensem de maneira inovadora e ambiciosa sobre algumas das seguintes questões : Qual é o problema?
Como resolver?
Como equilibrar os interesses dos envolvidos da comunidade?
Como garantir \textbf{que} aqueles com poder o usem com responsabilidade?
O \textbf{que} os jovens podem sugerir para provocar mudanças em suas comunidades e no entorno? ( HARVARD, 2020 ).
Como exemplo, Martins, Costa e Palhares ( 2018 ) apresentam \textbf{uma} atividade do Modelo ONU com estudantes de ensino médio sobre temas relacionados aos Direitos Humanos.
Os resultados permitiram desenvolver \textbf{uma} consciência sobre a cidadania global, assim como potencializaram o desenvolvimento de competências e habilidades nos desafios para a efetivação dos Direitos Humanos.
Outra prática pedagógica pode ser observada em Santos e Girotto ( 2017 ), numa atividade \textbf{que} proporcionou aos(às) estudantes assumir posições acerca de questões fundamentais das relações internacionais \textbf{contemporâneas}, como território, globalização, terrorismo, refugiados, entre outros, mediando situações-problema, nas quais se observou a modificação da relação dos/as estudantes com os conteúdos da temática.
O Modelo ONU permitiu a construção do \textbf{raciocínio} científico nas diversas escalas dos fenômenos trabalhados, problematizando e ampliando sua leitura a respeito da diversidade e da situação humana em diversos contextos. b ) Projeto Internacional “ Nós propomos : Geografia, Educação e Cidadania ” O projeto “ Nós propomos : Geografia, Educação e Cidadania ” foi criado em 2011/12 pelo prof. dr.
Sérgio Claudino, do Instituto de Geografia e Ordenamento do Território da Universidade de Lisboa/IGOT-ULisboa, Portugal.
Sua estratégia é a prática da Metodologia Ativa, com base nos estudos de caso, através dos quais os(as) estudantes devem pesquisar, identificar os problemas sociais e/ou ambientais locais, para \textbf{que} apresentem propostas concretas de resolução, partilhando com a comunidade.
Desde a sua iniciativa em Lisboa, expandiu -se para além de Portugal, para as regiões autónomas dos Açores e da Madeira, para o Brasil em 2014, a Espanha ( 2016/17 ), Moçambique ( 2017/18 ), Colômbia e Peru ( 2018 ) e México ( 2018/19 ).
Este projeto segue algumas etapas e protocolos \textbf{que} podem ser encontrados em Claudino e Coscurão ( 2019 ). \textbf{Uma} das vantagens positivas desta proposta \textbf{metodológica} é o engajamento dos(as) estudantes e professores junto à comunidade.
Os resultados podem ser apresentados junto à municipalidade através de encontros ou seminários.
Observa -se \textbf{que} as práticas escolares se tornam \textbf{um} palco em \textbf{que} os temas locais, regionais e mundiais são contextualizados, e neles os(as) estudantes se tornam protagonistas de seus espaços \textbf{imediatos}.
Escrevem Nascimento, Sartório e Claudino : > No âmbito local, os(as) estudantes podem perceber as suas conexões com as temáticas, e também estabelecer alianças transformadoras nas suas comunidades para ambientes melhores.
Pois, na escola, analisar o seu ambiente é desenvolver o senso crítico, permitir atitudes emancipatórias para mudanças de comportamento, incentivando o respeito à vida e disseminando novas ações ( 2019, p. 486 ).
Isto vale principalmente nas relações humanas e ambientais. c ) Sociodrama como prática escolar Jacob Levy Moreno ( 1889-1974 ), romeno-judeu e médico humanista, criou as terminologias psicodrama e sociodrama, no intuito de estimular as pessoas a descobrirem novas formas de estar no mundo, através da transformação social.
No sociodrama, utiliza -se a dramatização de situações.
Através dela, estabelecem -se papéis entre os participantes, com o objetivo de reconhecer os conflitos e buscar por soluções, fazendo com \textbf{que} encontrem outras formas de atuação social.
Os papéis sociais são maneiras \textbf{que} os \textbf{sujeitos} criam para se relacionar com todos e o mundo, permitindo diversas potencialidades ( MORENO, 1975 ).
Os papéis sociodramáticos são aqueles \textbf{que} representam ideias e experiências coletivas, identificando o grupo ou a coletividade \textbf{que} representa \textbf{uma} cultura.
A atuação do grupo é livre, lúdica, teatral e em \textbf{um} ambiente preparado para a prática.
Moreno ( 1994 ) destaca o teatro como linguagem espontânea para melhor representar os problemas sociais, e a relação interpessoal, por improvisos.
Cada \textbf{um} cria seu papel e texto a partir do tema apresentado, sendo o mais espontâneo possível, permitindo mudanças.
Como exemplo, Ferreira ( 2010 ) desenvolveu sua pesquisa de mestrado sobre a “ Expressividade do Sociodrama ” com estudantes do ensino médio e professores de todos os componentes curriculares, em \textbf{uma} escola de Taguatinga, Brasília, em 2010.
Para os(as) estudantes, a temática foi de relacionamentos, como bullying e outros contextos.
Para os professores foi a questão do prazer e sofrimento da prática docente. \textbf{Uma} das conclusões foi \textbf{que} a metodologia do sociodrama confirmou a expressividade para \textbf{uma} convivência saudável do grupo, amenizando os dramas e promovendo a espontaneidade e a criatividade.
Para Nery e Gisler ( 2019 ), os participantes do sociodrama são cocriadores e produtores de conhecimentos \textbf{que} contemplam histórias, crenças, afetos.
Manifestam -se pela fala, por comportamentos, posturas e dinâmicas grupais.
Os resultados vão além da análise da \textbf{cena}, de diálogos, personagens ; ultrapassam a observação dos comportamentos e atuações no campo social, grupal e dramático.
Avaliação No \textbf{que} se refere à avaliação, a elaboração do projeto político-pedagógico ( PPP ), além de \textbf{explicitar} a obrigatoriedade de desenvolver a temática “ História e Cultura Afro-Brasileira e Indígena ”, deve prever a concepção de avaliação do processo de ensino e aprendizagem em consonância com as legislações vigentes em âmbito federal e estadual.
Os critérios avaliativos devem levar em conta os objetivos propostos e a troca de conhecimento entre as diferentes áreas.
Assim, as contribuições \textbf{teóricas} advindas de todas as áreas do conhecimento precisam estar \textbf{explícitas}, para \textbf{que} os(as) estudantes compreendam a rede complexa do conhecimento ao longo da trilha integrada.
Além disso, é importante apresentar e debater/elucidar com eles(elas) o processo avaliativo, na perspectiva de \textbf{uma} avaliação processual, formativa e participativa.
Vale ressaltar \textbf{que} a Proposta Curricular de Santa Catarina orienta no sentido de \textbf{que} a avaliação ocorra de forma “ contínua, sistemática, expressa num movimento permanente de reflexão e ação ” ( \textbf{SANTA} CATARINA, 2014, p. 150 ).
A avaliação deve considerar processos de ensino e aprendizagem \textbf{que} enfatizem a produção de diferentes gêneros do discurso, leituras, desenvolvimento da oralidade em debates e seminários, pesquisa e elaboração de questões-problema, relatórios, textos argumentativos e dissertativos, projetos de pesquisa, criação de roteiros, produção audiovisual, confecção de mapas ( mentais ou \textbf{conceituais} ), propostas e projetos de intervenção na comunidade.
Também considera a capacidade de trabalhar em equipe, o relacionamento interpessoal e o desenvolvimento de observatórios, entre outras possibilidades a serem discutidas entre o coletivo da escola.
A orientação é \textbf{que} os aspectos qualitativos se sobreponham aos quantitativos e \textbf{que} os resultados da avaliação sejam indicadores de avanços, progressos e de crescimento para cada estudante, bem como apontem dificuldades e deficiências com vistas ao aprimoramento do processo educativo.

% matched lemmas: abordagem, agregar, aproximar, autor, cena, central, conceitual, configurar, conseguir, constituinte, constitutivo, contemporâneo, demandar, depender, desdobramento, despertar, dizer, estímulo, exclusão, explicitar, explícito, fomentar, freire, historicamente, histórico-cultural, hoje, humanos, imediato, inerente, influência, intermédio, legítimo, literário, metodológico, percepção, permear, popular, posicionamento, que, raciocínio, raça, remeter, ribeirinho, salientar, santo, singularidade, sistematizar, sistematização, sujeito, territorialidade, teórico, um, urgente, ver
\end{textsample}
